\scanpage{73}

% L01 -- top, bracketed and cancelled
\deleted{\uline{Choisissons} \inserted{\unsure{loc (Zar)}} section $s_0$ de $X$ admissible pour le syst. des $(\unsure{\tau_i^X|G^X|\xi_i^x})$ (déterminé mod homothétie dans $E=X\setminus C$ muni de \uline{l'origine $c$})}

% L02 -- upper third, crossed block
\deleted{(car $c\notin C$ en tt pt \inserted{i.e. $\rho$ inv} ($n$ impair et $\rho'$ inv., \deleted{\unsure{car}} ou $n$ pair \inserted{(alors $\rho=2\rho'$)} et $2$ inv)). -- \deleted{donc est} Donc on est dans les conditions suivantes (\unsure{identifions|identifier|identifiant} $E$ à $V$ grâce à l'origine $c=0$)}

% L03 -- upper third, item a
\deleted{\[
\text{a)}\quad \varphi_G^V:G\longrightarrow
\deleted{\unsure{\mathrm{Aut}(V,f)}}\mathrm O(f)
\]}

% L04 -- upper third, item b
\deleted{b) \deleted{\unsure{section $s_0$ de}} base $s_0$ de
$D_0^*=\bigcap_{1\le i\le n}H_i^{!}$
$(H_i^!=H_i\setminus H_i\cap\overline C)$}

% L05 -- left, beside item b
\margin{$a_{-1}^{\deleted{\unsure{*}}}$}

% L06 -- centre, introduction to three rows
On a \fbox{dans $V^\vee$}

% L07 -- centre, three rows
\[
\begin{aligned}
a_0^\vee&=\deleted{\unsure{\lambda_0-1}}(\lambda_0-1)a_0^*+\beta_0 a_1^*\\
a_i^\vee&=a_{i-1}^*+(\lambda_i-1)a_i^*+\beta_i a_{i+1}^*
 &1\le i\le n-1\\
a_n^\vee&=a_{n-1}^*+(\lambda_n-1)a_n^*
\end{aligned}
\]

% L08 -- left of the three rows, diagonal
% slanted
\margin{\uline{NB} $\rho$ inv. \unsure{exprime} ici que, posant $D_0=\bigcap_{1\le i\le n}H_i$, on a $a_0^\vee|D_0\ne0$ en chaque fibre (conséquence de la deuxième condition dans $2^\circ$)}

% L09 -- middle, cancelled line
\deleted{\unsure{Cela} \unsure{pour} $a_{-1}=$}

% L10 -- middle, warning in parentheses
(Attention au chgt des notations, il n'y a plus de $a_{-1},a_{-1}^*\ \ldots$)

% L11 -- below middle, starred calculation
\[
\deleted{\unsure{a'}}\quad\text{\textcircled{$*$}}\qquad
s_0\deleted{\unsure{=0-\frac1{\rho}(\zeta_0a_0+\zeta_1}}
\inserted{=\widetilde{s}_0=}
-\frac1\rho(\zeta_0a_0+\zeta_1a_1+\cdots+\zeta_n a_n)
\]

% L12 -- left, alongside coefficients
% slanted
\margin{On suppose pour simplifier maintenant les $\lambda_i=-1$ ($\tau_i$ \uline{réflexions})}

% L13 -- lower middle left
\[
\begin{aligned}
\zeta_n&=1\\
\zeta_{n-1}&=2\\
\zeta_i&=L_{n-i}(\beta_{n-1},\ldots,\beta_{i+1})
\end{aligned}
\]

% L14 -- lower middle right, square brackets
\[
\left[\begin{aligned}
\tau_i^V(s_0)&=s_0 &1\le\unsure{j|i}\le n\\
\text{i.e. }\langle s_0,a_i^\vee\rangle&=0 &1\le i\le n\\
\tau_0^V s_0&=s_0+a_0 &\text{i.e. }\langle s_0,a_0^\vee\rangle=1
\end{aligned}\right]
\]

% L15 -- below coefficients
\[
\zeta_0=L_n(\beta_{n-1},\ldots,\beta_1)=L_n(\beta_1,\ldots,\beta_{\unsure{n|n-1}})=\sigma
\]

% L16 -- bottom, item C
C) \unsure{Forme|Formes} quadratiques invariantes : \deleted{\unsure{ce sont les}} \inserted{elles forment} un \deleted{\unsure{torseur ss|torseur}} $\unsure{\mathcal O_S|G}$-Module \inserted{loc} \deleted{libre de} rg 1. On normalise par

% L17 -- bottom, left brace
\[
\left\{\begin{aligned}
f_V(a_0)&=1\quad\Longrightarrow\quad f_V(a_i)=\beta_0\beta_1\cdots\beta_{i-1}\\
\varphi(a_i,a_j)&=\deleted{\unsure{\{} }0\quad\text{si } i,j\text{ non consécutifs}\\
\varphi(a_i,a_{i+1})&=-\beta_0\cdots\beta_i
\end{aligned}\right.
\]
% archive sheet number: 72
